\section{Background and Related Work}
\label{sec:background}

\subsection{Constant-Time Programming}

A cryptographic implementation is \emph{constant-time} (CT) when execution is
independent of secret data. The usual leak surfaces are secret-dependent
branches, secret-dependent memory addresses, and secret-dependent
variable-latency instructions such as integer division~\cite{kocher}. Defenses
replace these with arithmetic selection, full scans or bit-slicing, and no
secrets in variable-latency instructions. Because compilers can reintroduce or
relocate a leak, CT checking must track the shipped build, not just the source.

\subsection{The Constant-Time Verification Landscape}

\paragraph{Dynamic taint analysis.}
ctgrind~\cite{ctgrind} marks secrets as uninitialized under Valgrind's
Memcheck~\cite{valgrind}, flagging tainted data in a branch or address -- a
precise binary-level check for the first two classes. Its blind spot is
variable-latency arithmetic.

\paragraph{Statistical black-box timing.}
dudect~\cite{dudect} applies a Welch $t$-test to cycle counts for fixed versus
random inputs; a difference signals a leak without modeling why. It is
complementary but noise-sensitive, and it localizes no cause.

\paragraph{Formal and static analysis.}
ct-verif~\cite{ctverif} verifies CT at the LLVM level by reduction to safety,
and Binsec/Rel~\cite{binsecrel} does relational symbolic execution on binaries.
They give stronger guarantees but require more setup and still encode a leak
model. No single technique covers branches, addresses, variable-latency
instructions, and black-box timing at once.

\subsection{Constant-Time Concerns in Post-Quantum Cryptography}

Kyber and Dilithium are lattice schemes~\cite{kyber,dilithium}. NIST standardized
them as ML-KEM and ML-DSA~\cite{fips203,fips204}. Their PQClean
implementations~\cite{pqclean} contain the constructs CT warns about.

KyberSlash~\cite{kyberslash} illustrates this: a division by a public modulus
compiles to a variable-time division on a secret-derived operand. Control flow
and addresses stay constant, so a branch-and-address check reports nothing; the
leak lives in the variable-latency class a Valgrind-style check cannot see.

Yet not every data-dependent timing is a defect. Dilithium signing uses rejection
sampling, so iteration count varies by design; FIPS~204 prescribes this, and the
variation is not a long-term-key leak. Distinguishing accepted behavior from a
genuine leak requires triage, not binary pass/fail.

\subsection{The Gap We Address}

The individual phenomena are known: KyberSlash division, build-sensitive CT, and
dudect timing. Missing is applying complementary checkers under one configuration
and build matrix, then reducing their output through default-deny triage. The
pipeline is Figure~\ref{fig:pipeline}; coverage and verdict classes are
Tables~\ref{tab:coverage} and~\ref{tab:corpus}.
